\chapter{Discussion}\label{ch:discussion}

\section{Answers to the research questions}

\paragraph{RQ1: How can data analytics improve business decision-making?}
The framework shows improvements at all three levels of analytics. At the \emph{descriptive} level, it replaces
manual spreadsheet consolidation with validated, reconciled KPIs that update with one command. At the
\emph{diagnostic} level, it reveals relationships that simple reports would miss. Growth is mostly price- and
mix-led. Discounts are not associated with larger orders. Order size depends on the channel, not the customer
type. At the \emph{predictive} level, it provides product-level forecasts with quantified uncertainty. Each of
these changes a concrete decision: discount budgets, sales-force allocation, production plans and logistics
capacity.

\paragraph{RQ2: Which preprocessing techniques are most effective?}
Three techniques proved most valuable. The first is \emph{deriving} missing values from accounting identities,
which gives exact restoration and should always be tried before statistical imputation. The second is
\emph{robust, group-wise outlier detection} (a per-product IQR on the log scale), which respects skewed
distributions and different product scales. Combined with the policy of flagging rather than deleting unusual
values, it keeps genuine bulk orders. The third is \emph{leakage-safe feature engineering}, in which history
features are shifted before rolling and every learned parameter is fitted on training data only. Robust
methods were preferable to normal-theory methods throughout, because the business variables are highly skewed.

\paragraph{RQ3: Which algorithm performs best?}
No single algorithm is best for every problem, and the margins between good models are small. For one-month-ahead
product demand, \BestReg was most accurate. \BestRec was most robust in the multi-step setting. At the company
level, the classical Holt--Winters model beat all machine-learning alternatives. For classification, logistic
regression matched the tree ensembles. This pattern agrees with \citet{makridakis2018}. Good features, a sound
validation design and benchmarking matter more than the choice of algorithm. A combined approach, with
statistical totals reconciled to machine-learned product mixes, is the most useful in practice. The
algorithm-level study in \cref{ch:algorithms} reached the same conclusion with independent code and a simpler
feature set. There, \AlgRegBest had the lowest monthly error (MAE \AlgRegBestMAE units), but linear regression
(\AlgLinMAE) and the random forest (\AlgRFMAE) were within a few units of it, and the seven classifiers spanned
only \AlgClsWorstAUC to \AlgClsBestAUC in ROC-AUC. Because the exact ranking of the leading models changes with
the features and tuning, the choice among them should favour simplicity and explainability.

\paragraph{RQ4: What patterns exist in the historical data?}
Strong year-end seasonality peaks in \PeakMonth. Growth is steady but price-driven, and demand is
price-inelastic. Revenue is geographically concentrated in Dhaka and Chattogram, and the portfolio is
concentrated in about ten products. Order size differs strongly by channel. Discounts have no detectable volume
effect, and weekday effects are absent. Segmentation adds a portfolio view (premium and gift packs, core family
packs, value single-serves and niche lines) and a market view (core retail, bulk distribution, online and
regional, emerging high-growth).

\paragraph{RQ5: How accurately can the future be predicted?}
Monthly demand per product can be forecast with a weighted error of about \BestRegWAPE\% one month ahead and
about \BestRecWAPE\% up to 21 months ahead. Company-level monthly units can be forecast with an error of about
\HWMAPE\%. Whether an individual order will be large is only moderately predictable (ROC-AUC \ClsAUC). The
limit lies in the information available, not in the models.

\paragraph{RQ6: How should results be presented to managers?}
Through an interactive dashboard organised around the management questions. It should have global filters,
KPI cards and simple, sorted charts; it should always show uncertainty for forecasts; and it should state
insights in plain language, each with an action. Publishing the dashboard and a summary website on free
platforms removes cost and installation barriers. In this project both are live (\cref{tab:links}), every page
has its own shareable address, and continuous deployment means that managers always see the latest results
without any action on their part.

\paragraph{RQ7: How much does predictive analytics add to historical reporting?}
A historical report can at best repeat last year's pattern, which is the seasonal-naive rule. The
machine-learning forecast reduces its error by \GainVsSNaive\%. Put differently, a planner using last year's
figures would make errors about \SNaiveWAPE\% of volume in each product-month, compared with \BestRegWAPE\% for
the framework. Prediction intervals add information that historical reporting cannot provide at all: how much
safety stock is needed.

\section{Managerial implications}
\begin{itemize}
  \item \textbf{Pricing and promotions.} Demand is inelastic and discounts do not raise order size. The discount
        budget of about \BDT\DiscountCost million over the period should be reviewed. Controlled experiments
        (A/B promotions) should establish whether targeted promotions pay off.
  \item \textbf{Volume growth.} Because growth relies on price, volume should become a separately tracked KPI,
        with growth initiatives in under-penetrated regions and in E-commerce.
  \item \textbf{Seasonal planning.} Production and distributor stock should be built from October for the
        \PeakMonth peak. February is the natural window for maintenance and promotions.
  \item \textbf{Portfolio.} The A-class products deserve the highest service levels. C-class products are
        candidates for review.
  \item \textbf{Logistics.} Large orders are concentrated in the distributor channel and in particular regions.
        Vehicle and stock allocation can be planned at that level.
  \item \textbf{Control.} A daily anomaly screen improves data quality and detects operational problems early.
\end{itemize}

\section{Limitations}\label{sec:limitations}
\begin{itemize}
  \item \textbf{Data provenance and realism.} The dataset is very clean: it has no missing values, its
        accounting identities hold exactly, and order size is almost independent of customer type, discount and
        product. Real operational data are usually messier. The pipeline was therefore tested on deliberately
        damaged copies, but conclusions about customer behaviour should be validated on further company data
        before they are acted on. The exact source and collection process of the dataset should be documented
        alongside this report.
  \item \textbf{Missing drivers.} No data were available on marketing spend, competitor prices, stock-outs,
        weather or holidays. Including them would probably improve both forecasting and classification.
  \item \textbf{Observational elasticity.} The elasticity estimate rests on list-price variation over time, not
        on randomised experiments. Prices changed by small amounts, so elasticity outside the observed range is
        unknown.
  \item \textbf{Future conditions.} All forecasts assume that past patterns continue. Structural breaks, such as
        a new competitor or an economic shock, would reduce accuracy.
  \item \textbf{Scope of evaluation.} The dashboard has not yet been evaluated with real managers. Deep-learning
        models were not explored, because the panel (52 series of 81 months) is too small to justify them.
\end{itemize}

\section{Reflection on the methodology}
Strict time-ordered evaluation with benchmarks was the most important methodological decision. Random
cross-validation would have produced flattering but misleading accuracy figures. Generating every reported
number from code turned out to be valuable as well: analyses were re-run many times as the work progressed, and
the report always remained consistent. Running a second, independent modelling study (\cref{ch:algorithms})
proved a cheap and effective way to test whether the conclusions depended on one implementation. Both practices
are recommended for future projects of this kind.
